\chapter{Algebraic Effect Handlers}
\label{chap:effects}

Algebraic Effects is the mathematical model for computational effects first introduced in \cite{plotkin2003algebraic} as an extension to the $\lambda$-calculus. After some time, Effect Handlers are introduced in \cite{plotkin2013handling} as a new addition to programming languages. Algebraic Effects and Handlers allows programmers to write programs with complex control flows. Functionalities such as exception handling, asynchronousity, I/O operations, and others are proven to be implement-able using Algebraic Effect Handlers. In this thesis, we use the term Algebraic Effect Handlers, Effect Handlers, and Handlers interchangably.

Algebraic Effect Handlers adapts the use of Call-By-Push-Value \cite{levy2013call} to separate value types, and computation types. This separation helps typing the computation together with possible effects, whereas value types are pure. This formal model usually provides reduction rules for its semantics.

In another spectrum, Effect Handlers are modeled by a combination of row types on top of \systemfw, named \systemfe \cite{xie2020effect}. The type rules are adjusted to accompany with a list of effects for an expression. This formal model usually provides evaluation contexts for its semantics.

\section{Examples}

\section{Language \lange}

In this section, we introduce \lange, a formal model of the Algebraic Effect Handlers as a model language for our target language for the compiler. Previously mentioned, we have two main choices for the formal language, (i) in the flavour of fine-grain call-by-value or (ii) in the flavour of \systemfe. We have chosen to adapt the (ii) flavour into the formal language due to its evaluation context semantics, which is helpful for our proofs.

Following \cite{xie2020effect}, we adopt the \textit{dot notation} to decompose and compose evaluation contexts, as defined in Figure \ref{fig:f-epsilon-dot}. We would write $v \cdot \text{handle } h \cdot E \cdot e$ as a shorthand for $v \, (\text{handle } h \, (E[e]))$.

\begin{singlespace}
\begin{figure}[h]
  \centering
  \begin{itemize}
      \item[] \begin{tabular}{llll}
          Context Composition & $E \cdot e$ & $\doteq$ & $E[e]$ \\
          Argument Evaluation & $\square \ e \cdot E$ & $\doteq$ & $E \ e$ \\
          Function Evaluation & $v \ \square \cdot E$ & $\doteq$ & $v \cdot E \ \doteq \ v \ E$ \\
          Handler Application & $\text{handle} \ h \ \square \cdot E$ & $\doteq$ & $\text{handle} \ h \cdot E \ \doteq \text{handle} \ h \ E$ 
      \end{tabular}
  \end{itemize}
  
  \caption{System $F^\epsilon$ Dot Notation}
  \label{fig:f-epsilon-dot}
\end{figure}

\end{singlespace}

\begin{singlespace}

\begin{figure}[h]
  \centering
  \begin{itemize}
      \item[] \begin{tabular}{llll}
          Expressions & $e$ & $::=$ & $v \mid e_1 \ e_2 \mid e[\sigma] \mid \text{handle}^\epsilon \ h \ e$ \\
          Values      & $v$ & $::=$ & $x \mid \lambda^\epsilon x:\sigma. e \mid \Lambda \alpha^k. v \mid \text{handler}^\epsilon \ h \mid \text{perform}^\epsilon \ op \ \overline{\sigma}$ \\
          Handlers    & $h$ & $::=$ & $\{op_1 \to f_1, \dots, op_n \to f_n\}$ \\
          Eval Contexts & $F$ & $::=$ & $\square \mid F \ e \mid v \ F \mid F[\sigma]$ \\
          & $E$ & $::=$ & $\square \mid E \ e \mid v \ E \mid E[\sigma] \mid \text{handle}^\epsilon \ h \ E$ \\
          Types       & $\sigma$ & $::=$ & $\alpha^k \mid c^k \ \sigma \dots \sigma \mid \sigma \to \epsilon \ \sigma \mid \forall \alpha^k. \sigma$ \\
          Kinds       & $k$ & $::=$ & $* \mid k \to k \mid \text{eff} \mid \text{lab}$
      \end{tabular}
  \end{itemize}
  \caption{System $F^\epsilon$ Syntax}
  \label{fig:f-epsilon-syntax}
\end{figure}

\begin{figure}[h]
  \centering
  \begin{mathpar}
  \inferrule[T-Var]
    {x : \sigma \in \Gamma}
    {\Gamma \vdash_{\text{val}} x : \sigma}

  \inferrule[T-Abs]
    {\Gamma, x : \sigma_1 \vdash e : \sigma_2 \mid \epsilon}
    {\Gamma \vdash_{\text{val}} \lambda^\epsilon x : \sigma_1 . e : \sigma_1 \to^\epsilon \sigma_2}

  \inferrule[T-Val]
    {\Gamma \vdash_{\text{val}} v : \sigma}
    {\Gamma; w \vdash v : \sigma \mid \epsilon}

  \inferrule[T-App]
    {\Gamma; w \vdash e_1 : \sigma_1 \to^\epsilon \sigma \mid \epsilon \\ \Gamma; w \vdash e_2 : \sigma_1 \mid \epsilon}
    {\Gamma; w \vdash e_1 \, e_2 : \sigma \mid \epsilon}

  \inferrule[T-Handle]
    {\Gamma \vdash_{\text{ops}} h : \sigma \mid l \mid \epsilon \\ \Gamma; w \vdash e : \sigma \mid \langle l \mid \epsilon \rangle}
    {\Gamma; w \vdash \text{handle}^\epsilon \, h \, e : \sigma \mid \epsilon}
  \end{mathpar}
  \caption{System $F^\epsilon$ Selected Type Rules}
  \label{fig:f-epsilon-type}
\end{figure}

\begin{figure}[h]
  \centering
  \begin{mathpar}
  \inferrule[App]
    {\strut}
    {(\lambda^\epsilon x : \sigma . e) \, v \longrightarrow e[x := v]}

  \inferrule[Handler]
    {\strut}
    {(\text{handler}^\epsilon \, h) \, v \longrightarrow \text{handle}^\epsilon \, h \cdot v \, ()}

  \inferrule[Return]
    {\strut}
    {\text{handle}^\epsilon \, h \cdot v \longrightarrow v}

  \inferrule[Perform]
    {op \notin \text{bop}(E) \wedge (op \to f) \in h \\
     op : \forall \overline{\alpha} . \sigma_1 \to \sigma_2 \in \Sigma(l) \\
     k = \lambda^\epsilon x : \sigma_2[\overline{\alpha} := \overline{\sigma}] . \text{handle}^\epsilon \, h \cdot E \cdot x}
    {\text{handle}^\epsilon \, h \cdot E \cdot \text{perform} \, op \, \overline{\sigma} \, v \longrightarrow f[\overline{\sigma}] \, v \, k}
  \end{mathpar}
  \caption{System $F^\epsilon$ Reduction Rules}
  \label{fig:f-epsilon-reduction}
\end{figure}

\end{singlespace}
